%========================================
% MatinBook - Stage 13: Index Test
% Test: makeindex + Index Entries + Subentries + Cross-references
% Purpose: Validate index generation system
%========================================

\makeatletter
\def\input@path{{../}}
\makeatother

\documentclass{matinbook}

\title{آزمون سیستم نمایه}
\author{تیم توسعه MatinBook}
\date{تابستان ۱۴۰۵}

\begin{document}

%----------------------------------------
% Title Page
%----------------------------------------
\maketitle

%----------------------------------------
% Chapter 1: Programming Concepts
%----------------------------------------
\chapter{مفاهیم برنامه‌نویسی}

\section{زبان‌های برنامه‌نویسی}

\subsection{پایتون}

پایتون\index{پایتون} یک زبان برنامه‌نویسی
سطح بالا، تفسیری و همه‌منظوره است.
این زبان توسط خیدو فان روسوم\index{فان روسوم، خیدو}
در سال ۱۹۹۱ ساخته شد.

ویژگی‌های اصلی پایتون\index{پایتون!ویژگی‌ها}
عبارتند از:

\begin{itemize}
    \item \textbf{خوانایی بالا:} نحو ساده و شبیه به
    زبان طبیعی\index{پایتون!نحو}
    \item \textbf{کتابخانه غنی:} کتابخانه استاندارد
    گسترده\index{پایتون!کتابخانه استاندارد}
    \item \textbf{چندسکویی:} اجرا روی سیستم‌عامل‌های
    مختلف\index{چندسکویی}
\end{itemize}

\subsection{زبان C++}

C++\index{C++} یک زبان برنامه‌نویسی
سطح میانی با قابلیت‌های شیءگرایی\index{شیءگرایی}
است. این زبان توسط بیارنه استراستروپ\index{استراستروپ، بیارنه}
در سال ۱۹۸۵ توسعه یافت.

ویژگی‌های C++\index{C++!ویژگی‌ها}:

\begin{itemize}
    \item \textbf{سرعت بالا:} کامپایل مستقیم به کد ماشین
    \index{C++!سرعت}
    \item \textbf{مدیریت حافظه:} کنترل دستی حافظه
    \index{مدیریت حافظه}
    \item \textbf{برنامه‌نویسی جنریک:} قالب‌ها
    \index{برنامه‌نویسی جنریک}
\end{itemize}

\section{الگوریتم‌ها}

الگوریتم\index{الگوریتم} مجموعه‌ای از دستورالعمل‌های
گام‌به‌گام برای حل یک مسئله است.

\subsection{انواع الگوریتم‌ها}

\begin{itemize}
    \item \textbf{جستجو:}
    جستجوی خطی\index{الگوریتم!جستجوی خطی} و
    جستجوی دودویی\index{الگوریتم!جستجوی دودویی}
    از الگوریتم‌های پایه جستجو هستند.
    
    \item \textbf{مرتب‌سازی:}
    مرتب‌سازی سریع\index{الگوریتم!مرتب‌سازی سریع} و
    مرتب‌سازی ادغامی\index{الگوریتم!مرتب‌سازی ادغامی}
    از پرکاربردترین الگوریتم‌های مرتب‌سازی هستند.
    
    \item \textbf{گراف:}
    الگوریتم دایکسترا\index{الگوریتم!دایکسترا} برای
    یافتن کوتاه‌ترین مسیر\index{کوتاه‌ترین مسیر}
    در گراف‌های وزن‌دار\index{گراف!وزن‌دار} استفاده می‌شود.
\end{itemize}

\section{ساختمان داده‌ها}

ساختمان داده\index{ساختمان داده} روشی برای
سازمان‌دهی و ذخیره‌سازی داده‌ها در کامپیوتر است.

انواع ساختمان داده\index{ساختمان داده!انواع}:

\begin{itemize}
    \item \textbf{آرایه}\index{آرایه}:
    مجموعه‌ای از عناصر با اندیس\index{اندیس} عددی
    \item \textbf{لیست پیوندی}\index{لیست پیوندی}:
    مجموعه‌ای از گره‌ها\index{گره} که هر کدام
    به عنصر بعدی اشاره می‌کنند
    \item \textbf{پشته}\index{پشته}:
    ساختار \lr{LIFO}\index{LIFO}
    (آخرین ورودی، اولین خروجی)
    \item \textbf{صف}\index{صف}:
    ساختار \lr{FIFO}\index{FIFO}
    (اولین ورودی، اولین خروجی)
    \item \textbf{درخت}\index{درخت}:
    ساختار سلسله‌مراتبی\index{ساختار سلسله‌مراتبی}
    با ریشه\index{ریشه} و برگ‌ها\index{برگ}
    \item \textbf{گراف}\index{گراف}:
    مجموعه‌ای از رئوس\index{رأس} و یال‌ها\index{یال}
    \item \textbf{جدول درهم‌سازی}\index{جدول درهم‌سازی}:
    ساختار کلید-مقدار\index{کلید-مقدار} با دسترسی $O(1)$
\end{itemize}

%----------------------------------------
% Chapter 2: Mathematical Concepts
%----------------------------------------
\chapter{مفاهیم ریاضی}

\section{نظریه اعداد}

نظریه اعداد\index{نظریه اعداد} شاخه‌ای از
ریاضیات\index{ریاضیات} است که به مطالعه
اعداد صحیح\index{اعداد صحیح} می‌پردازد.

\subsection{اعداد اول}

عدد اول\index{عدد اول} عددی طبیعی بزرگتر از ۱
است که فقط بر ۱ و خودش بخش‌پذیر باشد.

قضیه اساسی حساب\index{قضیه!اساسی حساب}
بیان می‌کند که هر عدد طبیعی بزرگتر از ۱
به صورت حاصلضرب اعداد اول\index{اعداد اول!تجزیه}
قابل نمایش است.

\subsection{بزرگترین مقسوم‌علیه مشترک}

بزرگترین مقسوم‌علیه مشترک\index{GCD}
یا \lr{GCD}\index{GCD|seealso{بزرگترین مقسوم‌علیه مشترک}}
دو عدد، بزرگترین عددی است که هر دو را می‌شمارد.

الگوریتم اقلیدس\index{اقلیدس!الگوریتم}
روشی کارآمد برای محاسبه \lr{GCD} است.

\section{جبر}

جبر\index{جبر} شاخه‌ای از ریاضیات است
که به مطالعه ساختارهای جبری\index{ساختار جبری}
می‌پردازد.

\subsection{گروه}

گروه\index{گروه} یک ساختار جبری شامل
یک مجموعه و یک عمل دوتایی\index{عمل دوتایی}
با چهار ویژگی بسته بودن\index{گروه!بسته بودن}،
شرکت‌پذیری\index{گروه!شرکت‌پذیری}،
وجود عضو خنثی\index{گروه!عضو خنثی} و
وجود عضو وارون\index{گروه!عضو وارون} است.

\subsection{ماتریس}

ماتریس\index{ماتریس} آرایه‌ای مستطیلی
از اعداد است. ماتریس‌ها\index{ماتریس!انواع}
در جبر خطی\index{جبر خطی}،
گرافیک کامپیوتری\index{گرافیک کامپیوتری}
و یادگیری ماشین\index{یادگیری ماشین}
کاربرد دارند.

\section{آنالیز ریاضی}

آنالیز ریاضی\index{آنالیز ریاضی} به مطالعه
حد\index{حد}، پیوستگی\index{پیوستگی}،
مشتق\index{مشتق} و انتگرال\index{انتگرال}
می‌پردازد.

قضیه مقدار میانی\index{قضیه!مقدار میانی}
و قضیه اساسی حسابان\index{قضیه!اساسی حسابان}
از قضایای مهم آنالیز هستند.

%----------------------------------------
% Chapter 3: Using Index Commands
%----------------------------------------
\chapter{راهنمای استفاده از نمایه}

\section{دستورات پایه}

برای ایجاد مدخل در نمایه از دستور
\texttt{\textbackslash index} استفاده می‌شود:

\begin{itemize}
    \item \texttt{\textbackslash index\{واژه\}}
    --- مدخل اصلی\index{مدخل اصلی}
    \item \texttt{\textbackslash index\{واژه!زیرواژه\}}
    --- زیرمدخل\index{زیرمدخل}
    \item \texttt{\textbackslash index\{واژه!زیرواژه!زیرزیرواژه\}}
    --- زیرزیرمدخل\index{زیرزیرمدخل}
\end{itemize}

\section{دستورات پیشرفته}

\lr{MatinBook} دستورات کمکی زیر را برای نمایه
فراهم می‌کند:

\begin{itemize}
    \item \texttt{\textbackslash idxbold}:
    مدخل با فونت سیاه\index{سیاه|textbf}
    \item \texttt{\textbackslash idxsee}:
    ارجاع «نگاه کنید به»\index{ارجاع!نگاه کنید به}
    \item \texttt{\textbackslash idxseealso}:
    ارجاع «همچنین نگاه کنید به»\index{ارجاع!همچنین نگاه کنید به}
    \item \texttt{\textbackslash idxsub}:
    ایجاد زیرمدخل سریع
    \item \texttt{\textbackslash idxsubsub}:
    ایجاد زیرزیرمدخل سریع
\end{itemize}

\section{نمونه ارجاع‌های نمایه}

\begin{itemize}
    \item \idxsee{مرتب‌سازی سریع}{الگوریتم!مرتب‌سازی سریع}
    \item \idxseealso{درخت دودویی}{درخت}
    \item \idxbold{مفهوم مهم}
    \item \idxsub{برنامه‌نویسی}{پایتون}
\end{itemize}

%----------------------------------------
% Summary and Index
%----------------------------------------
\chapter{خلاصه نتایج}

\section{وضعیت سیستم نمایه}

\begin{table}[h]
\centering
\caption{وضعیت قابلیت‌های نمایه}
\label{tab:index-status}
\begin{tabular}{|c|C{6cm}|c|}
\hline
\textbf{قابلیت} & \textbf{توضیحات} & \textbf{وضعیت} \\
\hline
مدخل اصلی & \lr{\textbackslash index\{واژه\}} & \textcolor{matingreen}{✅} \\
زیرمدخل & \lr{\textbackslash index\{واژه!زیر\}} & \textcolor{matingreen}{✅} \\
زیرزیرمدخل & \lr{\textbackslash index\{واژه!زیر!زیر\}} & \textcolor{matingreen}{✅} \\
ارجاع \lr{see} & \lr{\textbackslash idxsee\{از\}\{به\}} & \textcolor{matingreen}{✅} \\
ارجاع \lr{see also} & \lr{\textbackslash idxseealso\{از\}\{به\}} & \textcolor{matingreen}{✅} \\
فونت سیاه & \lr{\textbackslash idxbold\{واژه\}} & \textcolor{matingreen}{✅} \\
زیرمدخل سریع & \lr{\textbackslash idxsub\{اصلی\}\{زیر\}} & \textcolor{matingreen}{✅} \\
چاپ نمایه & \lr{\textbackslash printindex} & \textcolor{matingreen}{✅} \\
\hline
\end{tabular}
\end{table}

\section{نتیجه}

\textbf{سیستم نمایه \lr{MatinBook} نسخه ۱ با موفقیت آزمایش شد.}
بیش از ۶۰ مدخل در ۲ فصل موضوعی با زیرمدخل‌ها
و ارجاع‌های متقابل به درستی در نمایه ثبت می‌شوند. 🎉

%----------------------------------------
% Print Index
%----------------------------------------
\printindex

\end{document}